\documentclass[11pt,a4paper]{article}
\usepackage[margin=2.5cm]{geometry}
\usepackage[hidelinks]{hyperref}
\usepackage{amsmath,amssymb,amsthm}
\usepackage{booktabs}
\usepackage{lineno}
\usepackage{microtype}

\newtheorem{proposition}{Proposition}
\newtheorem{conjecture}{Conjecture}
\theoremstyle{remark}
\newtheorem{remark}{Remark}

\newcommand{\PSL}{\mathrm{PSL}(2,\mathbb{C})}
\newcommand{\mfld}[1]{\texttt{#1}}
\newcommand{\word}[1]{\texttt{#1}}
\newcommand{\degr}{^{\circ}}
\newcommand{\Zfive}{\mathbb{Z}/5}
\linenumbers

\begin{document}

\title{Homology Class Asymmetry in the Loxodromic\\
Twist Spectrum of Compact Hyperbolic 3-Manifolds}

\author{M.~L.~Gentry%
\thanks{Independent Researcher, Seattle, WA, USA;
\texttt{drmlgentry@protonmail.com};
ORCID: 0009-0006-4550-2663}}
\date{\today}
\maketitle

%------------------------------------------------------------------
\begin{abstract}
We introduce a homology-resolved spectral floor invariant for compact
hyperbolic 3-manifolds with $H_1(M,\mathbb{Z})=\Zfive$.
The abelianization map partitions loxodromic elements into five
homology classes; the spectral floor
$\phi_{\rm floor}^{(k)}(L)$ is the minimum folded twist angle
among loxodromics of word length $\leq L$ in class $k$.
A complete census of $\mfld{m006}$ (OrientableClosedCensus[43],
vol $= 2.029$, $H_1=\Zfive$) to word length 12
($1{,}062{,}802$ genuine loxodromics) shows that all five floors
decay exponentially with word length, with class-dependent rates
forming a paired hierarchy:
$c_1=c_4\approx 0.985$, $c_2=c_3\approx 0.401$, $c_0\approx 0.750$,
confirmed by AIC model selection against power-law alternatives
($\Delta\mathrm{AIC}>40$ for classes 0, 1, 4).
At length 12 the floors satisfy an approximate $4{:}2{:}1$ ratio.
The same pairing structure appears for the Meyerhoff manifold
$\mfld{m003}$ (also $H_1=\Zfive$), suggesting the hierarchy is
a property of $\Zfive$ torsion rather than of either specific manifold.
We conjecture that the pairing $c_k=c_{5-k}$ reflects complex
conjugation in the invariant trace field.
\end{abstract}

%------------------------------------------------------------------
\section{Introduction}
\label{sec:intro}

The complex length spectrum
$\{\ell(\gamma)+i\phi(\gamma)\}_\gamma$ of a compact hyperbolic
3-manifold is a fundamental geometric invariant.
For manifolds with non-trivial first homology, the abelianization
$\pi_1(M)\to H_1(M,\mathbb{Z})$ partitions geodesics into homology
classes. A natural question is whether the twist angle distribution
depends on the homology class.

We study two compact arithmetic hyperbolic 3-manifolds with
$H_1=\Zfive$: the Meyerhoff manifold
$M_1 = \mfld{m003}$ (OrientableClosedCensus[1], vol $=0.9814$)
and $M_2 = \mfld{m006}$ (OrientableClosedCensus[43],
vol $=2.029$)~\cite{GMM,MaclachlanReid}.
Both arise in the Hyperbolic Flavor Geometry
programme~\cite{GentryUnified} as geometric models for
Standard Model lepton and quark mixing.

We define the \emph{spectral floor}
\begin{equation}
\phi_{\rm floor}^{(k)}(L)
= \min\bigl\{\phi_{\rm fold}(\gamma) :
  [\gamma]=k,\; |w(\gamma)|\leq L\bigr\},
\end{equation}
where $\phi_{\rm fold}=\min(|\phi|,\pi-|\phi|)\in[0,\pi/2]$
and $[w]= \sum_i\varepsilon(g_i)\bmod 5$
with $\varepsilon(g)=+1$ (lowercase) or $-1$ (uppercase).
We show that this invariant reveals a class-dependent hierarchy
absent from the unpartitioned spectrum.

%------------------------------------------------------------------
\section{Setup and census}
\label{sec:setup}

Twist angles $\phi(\gamma)=\arg\lambda(\gamma)$ are computed via
SnapPy's \texttt{G.SL2C(word)} at 150-bit precision~\cite{SnapPy}.
We enumerate all reduced words up to length $L$ over generators
$\{a,b,A,B\}$ ($A=a^{-1}$, $B=b^{-1}$), retain genuine
loxodromics ($|\mathrm{Tr}(\rho(w))|>2.01$, equivalently
$|\lambda|>1.01$), and record $\phi_{\rm fold}(w)$ per
homology class $[w]\in\Zfive$.
At length 12: $1{,}062{,}802$ genuine loxodromics from
$1{,}062{,}880$ words evaluated in $84.6$~seconds on a
standard laptop. Near-$\pi$ twists are verified via
$|\mathrm{Im}(\lambda)/\mathrm{Re}(\lambda)|$ directly.

%------------------------------------------------------------------
\section{Results}
\label{sec:results}

\subsection{The 4:2:1 hierarchy at length 12}

Table~\ref{tab:floors_L12} gives the length-12 spectral floors
per homology class for $\mfld{m006}$.
The five classes split into three groups with floors in ratio
approximately $4{:}2{:}1$:

\begin{table}[htbp]
\centering
\caption{Spectral floors per homology class at word length $\leq 12$
for $\mfld{m006}$ ($1{,}062{,}802$ loxodromics).
The ratio $\phi^{(0)}:\phi^{(2,3)}:\phi^{(1,4)}\approx 4{:}2{:}1$
is stable across lengths $\geq 10$.}
\label{tab:floors_L12}
\begin{tabular}{clll}
\toprule
Class $k$ & $\phi_{\rm floor}^{(k)}(12)$ (${}^\circ$) & Word & Length \\
\midrule
0 & 0.02300 & \word{aaaaaaaaabAb} & 12 \\
1 & 0.00530 & \word{aaaaaababAAA} & 12 \\
2 & 0.01060 & \word{aaababaaabab} & 12 \\
3 & 0.01060 & \word{AAABABAAABAB} & 12 \\
4 & 0.00530 & \word{aaaBAAABAAAA} & 12 \\
\midrule
\multicolumn{2}{l}{Ratio $\phi^{(0)}:\phi^{(2,3)}:\phi^{(1,4)}$}
& \multicolumn{2}{l}{$4.3:2:1$} \\
\bottomrule
\end{tabular}
\end{table}

The pairing $\phi^{(1)}=\phi^{(4)}$ and $\phi^{(2)}=\phi^{(3)}$
reflects the involution $k\mapsto -k\pmod{5}$ on $H_1=\Zfive$.

\subsection{Exponential decay with class-dependent rates}

Over the range $L=1,\ldots,12$, each class exhibits log-linear decay:
\begin{equation}
\phi_{\rm floor}^{(k)}(L) \sim A_k\,e^{-c_k L}.
\label{eq:exp_decay}
\end{equation}
Table~\ref{tab:rates} gives the fitted rates $c_k$ with bootstrap
confidence intervals and AIC model comparison.

\begin{table}[htbp]
\centering
\caption{Exponential decay rates $c_k$ from log-linear fits
to $\phi_{\rm floor}^{(k)}(L)$, $L=1,\ldots,12$.
Bootstrap 95\% CIs (2000 resamples).
$\Delta\mathrm{AIC} = \mathrm{AIC}(\text{exp}) - \mathrm{AIC}(\text{pow})$;
negative values favour exponential over power-law.
The CIs of $c_1,c_4$ overlap completely; likewise $c_2,c_3$.}
\label{tab:rates}
\begin{tabular}{cccr}
\toprule
Class $k$ & $c_k$ & 95\% CI & $\Delta$AIC \\
\midrule
0 & 0.750 & [0.434,\ 0.944] & $-33.6$ \\
1 & 0.985 & [0.541,\ 1.376] & $-49.0$ \\
2 & 0.401 & [0.080,\ 0.645] & $-4.0$  \\
3 & 0.401 & [0.079,\ 0.651] & $-4.0$  \\
4 & 0.985 & [0.585,\ 1.431] & $-49.0$ \\
\bottomrule
\end{tabular}
\end{table}

The exponential model is preferred for all five classes.
The pairing $c_1=c_4\approx 0.985$ (fastest decay) and
$c_2=c_3\approx 0.401$ (slowest) with class~0 intermediate
($c_0\approx 0.750$) is the principal structural result.

\subsection{Universality: Meyerhoff manifold}

Table~\ref{tab:both_L12} shows that $\mfld{m003}$
(Meyerhoff manifold) exhibits the \emph{same} pairing at length 12,
with ratio $\phi^{(2,3)}/\phi^{(1,4)}\approx 2.17$.

\begin{table}[htbp]
\centering
\caption{Length-12 spectral floors for both manifolds.
Both exhibit the pairing $\phi^{(1,4)}<\phi^{(0)}<\phi^{(2,3)}$
consistent with the involution $k\mapsto{-k}\pmod 5$.}
\label{tab:both_L12}
\begin{tabular}{cllll}
\toprule
Class $k$ & \multicolumn{2}{c}{$\mfld{m003}$}
           & \multicolumn{2}{c}{$\mfld{m006}$} \\
           & floor ($\degr$) & Word & floor ($\degr$) & Word \\
\midrule
0 & 0.04373 & \word{aaaabbbaabaB} & 0.02300 & \word{aaaaaaaaabAb} \\
1 & 0.04319 & \word{aaBAbAAbAABB} & 0.00530 & \word{aaaaaababAAA} \\
2 & 0.09387 & \word{aabAAbAAABB}  & 0.01060 & \word{aaababaaabab} \\
3 & 0.09387 & \word{aaaaBabbAAB}  & 0.01060 & \word{AAABABAAABAB} \\
4 & 0.04319 & \word{aabbAAbaBaaB} & 0.00530 & \word{aaaBAAABAAAA} \\
\midrule
\multicolumn{2}{l}{$\phi^{(2,3)}/\phi^{(1,4)}$} & $2.17$ & & $2.00$ \\
\bottomrule
\end{tabular}
\end{table}

The appearance of the same pairing in two arithmetically distinct
manifolds sharing only $H_1=\Zfive$ suggests the hierarchy is
a property of the $\Zfive$ torsion structure.

\subsection{The twist-suppressed geodesic}

The class-4 floor is set from length 6 by word \word{AAABAB}:
\begin{align*}
|\lambda(\word{AAABAB})| &= 3.730, &
\phi(\word{AAABAB}) &= 179.995\degr, &
\phi_{\rm fold} &= 0.0053\degr, \\
\ell(\word{AAABAB}) &= 2.633, &
|\mathrm{Tr}| &= 3.998,  &
[\word{AAABAB}] &= 4 \in \Zfive.
\end{align*}
The element is a near-pure translation ($\lambda\approx -3.730$,
nearly real negative) in a non-identity homology class,
demonstrating that twist suppression and non-trivial homology
class are geometrically compatible.

%------------------------------------------------------------------
\section{Discussion}
\label{sec:discussion}

\begin{proposition}
\label{prop:zero}
For any compact hyperbolic 3-manifold $M$ with $H_1=\Zfive$
and any class $k\in\Zfive$,
$\phi_{\rm floor}^{(k)}(L)\to 0$ as $L\to\infty$.
\end{proposition}

\begin{proof}
By the Ambient Prime Geodesic Theorem~\cite{BalogBiggs23} and
equidistribution of homology classes~\cite{Sunada85},
the count of class-$k$ geodesics with $\phi_{\rm fold}<\epsilon$
and length $\leq T$ grows as $(4\epsilon/\pi)\cdot e^{2T}/(10T)$,
which diverges for any fixed $\epsilon>0$.
\end{proof}

\begin{remark}
Proposition~\ref{prop:zero} guarantees convergence to zero but
not the rate. The class-dependent rates $c_k$ (Table~\ref{tab:rates})
are the new content of this paper.
The exponential fits are empirical over $L=1,\ldots,12$;
no claim is made regarding the asymptotic rate as $L\to\infty$.
\end{remark}

The pairing $c_1=c_4$, $c_2=c_3$ reflects the involution
$k\mapsto -k\pmod{5}$ on $\Zfive$, since $1+4\equiv 2+3\equiv 0$.
We propose that this is arithmetic in origin:

\begin{conjecture}
The rate pairing $c_k=c_{5-k}$ ($k\neq 0$) arises from the
outer involution $\rho\mapsto\bar{\rho}$ of the holonomy
representation induced by complex conjugation in the invariant
trace field of $\mfld{m006}$.
Manifolds sharing the same invariant trace field and quaternion
algebra have the same decay rate hierarchy.
\end{conjecture}

The rates $c_k$ constitute a new family of geometric invariants
of compact hyperbolic 3-manifolds with $H_1\cong\Zfive$.
Extending the census across the OrientableClosedCensus for
other manifolds with $H_1=\mathbb{Z}/p$ ($p$ prime) would test
whether the rate hierarchy depends on $p$, on arithmeticity,
or on the specific quaternion algebra.

%------------------------------------------------------------------
\section{Conclusions}
\label{sec:concl}

The homology-resolved spectral floor $\phi_{\rm floor}^{(k)}(L)$
reveals a class-dependent $4{:}2{:}1$ hierarchy in the loxodromic
twist spectra of two compact hyperbolic 3-manifolds with $H_1=\Zfive$.
Exponential decay with paired rates $c_1=c_4\approx 0.985$,
$c_2=c_3\approx 0.401$, $c_0\approx 0.750$ is confirmed by AIC
model selection over $L=1,\ldots,12$.
The hierarchy persists in both manifolds despite their differing
arithmetic invariants, implicating the $\Zfive$ torsion structure
as the origin.
A conjectured arithmetic explanation via complex conjugation in
the invariant trace field is proposed for future verification
via the Snap package.

\section*{Data Availability}
Complete census files and scripts at
\url{https://github.com/drmlgentry/hyperbolic-flavor-scan}.

\section*{Acknowledgements}
The author thanks the SnapPy development team.
During preparation of this manuscript the author used Claude
(claude-sonnet-4-6, Anthropic) as a research assistant.

\begin{thebibliography}{99}

\bibitem{GMM}
D.~Gabai, R.~Meyerhoff, P.~Milley,
J.\ Amer.\ Math.\ Soc.\ \textbf{22}, 1157 (2009).

\bibitem{MaclachlanReid}
C.~Maclachlan, A.~W.~Reid,
\emph{The Arithmetic of Hyperbolic 3-Manifolds}, Springer (2003).

\bibitem{GentryUnified}
M.~L.~Gentry,
``Hyperbolic Flavor Geometry,''
SSRN \href{https://doi.org/10.2139/ssrn.6775158}{6775158} (2026);
PTEP Paper No.\ 2605-035.

\bibitem{GentryCKM}
M.~L.~Gentry,
``Quark Mixing from Hyperbolic Geometry,''
SSRN \href{https://doi.org/10.2139/ssrn.6583550}{6583550} (2026).

\bibitem{GentryPMNS}
M.~L.~Gentry,
``Lepton Mixing from Borel Structure of Hyperbolic Holonomy,''
SSRN \href{https://doi.org/10.2139/ssrn.6583553}{6583553} (2026);
Phys.\ Lett.\ B PLB-D-26-01448.

\bibitem{BalogBiggs23}
A.~Balog et al.,
``Ambient Prime Geodesic Theorems on Hyperbolic 3-Manifolds,''
Int.\ Math.\ Res.\ Not.\ \textbf{2023}(1), 588--606.

\bibitem{Margulis69}
G.~A.~Margulis,
Funct.\ Anal.\ Appl.\ \textbf{3}, 335 (1969).

\bibitem{Sunada85}
T.~Sunada,
Ann.\ Math.\ \textbf{121}, 169 (1985).

\bibitem{Chinburg01}
T.~Chinburg, E.~Friedman, K.~N.~Jones, A.~W.~Reid,
Ann.\ Scuola Norm.\ Sup.\ Pisa \textbf{30}, 1 (2001).

\bibitem{SnapPy}
M.~Culler, N.~M.~Dunfield, M.~Goerner, J.~R.~Weeks,
SnapPy, \url{http://snappy.computop.org}.

\bibitem{Snap}
D.~Coulson, O.~Goodman, C.~Hodgson, W.~Neumann,
Experimental Math.\ \textbf{9}, 127 (2000).

\end{thebibliography}

\end{document}
